\section{Grundlagen und Stand der Forschung}\label{sec:Grundlagen}
Hier eine kurze Übersicht zu:
- BIM und IFC \newline
- Kurze Beschreibung der genutzten IFC Entitäten. \newline

Eingehen auf die That Open Bibliothek \newline
- Fragments zur Performance Verbesserung \newline
- IFC-Modell vs. Rendering-Modell \newline
- Kurz erklären, warum die ursprünglichen IFC-Quelldaten unabhängig von der Fragments transformation erhalten werden sollen. \newline


Mobile Navigation in Gebäuden. \newline
- Waypoints \newline
- LIDAR zur positionsbestimmung über hochauflösende Punktewolken \newline

Quellen: \newline
Improving autonomous robotic navigation using IFC files.pdf \newline
Behandelt fast perfekt die Schritte nach der Annotation der IFC-Datei: \newline
- Extraktion von Gebäudeinformationen \newline
- Navigationskarte und semantische Informationen \newline
- Wegpunkte \newline
- ROS- beziehungsweise Roboterausführung \newline


BIM-to-Robot\_Mapping\_Constructing\_IFC-Referenced\_Occupancy\_Grids\_and\_Semantic\_Metadata\_for\_Rule-Based\_Navigation.pdf \newline
Behandelt auch die Lücke zwischen Annotation und Roboter \newline
- 2D-Navigationskarte aus einem IFC-Modell \newline
- IFC-Semantik für Roboterlokalisierung und Pfadplanung \newline

Component-based robot prefabricated construction simulation using.pdf \newline
- Identität über GlobalId \newline
- Simulationsarchitektur evtl. interessant. \newline

Automated Robotic Planning for Corrective.pdf \newline
- Raumnavigation und Wegpunktgenerierung \newline
- LLM-Erweiterung \newline

BIM-ROS Integration.pdf (Kein Paper, Aber offizieller buildingSMART-Use-Case) \newline
- Integration von BIM und Robotik \newline
- IFC-basierte Datenbereitstellung für ROS \newline

